\chapter{Introduction}
\label{ch:intro}

Capacity planning asks how much computing capacity a service needs in order to maintain acceptable performance as load changes. In modern containerised systems this question is no longer answered reliably by static rules of thumb. A microservice may be deployed inside Docker, constrained by CPU quotas, scheduled by the host operating system, and executed by language runtimes with very different concurrency models. Under these conditions, response time can change sharply as utilisation grows, and the cause of the change is often not visible from latency graphs alone.

Classical queueing theory provides the conceptual starting point. In an M/G/$c$ model, arrivals are usually assumed to be Poisson, service times follow a fixed general distribution, and $c$ workers serve jobs from a queue. This abstraction is valuable because it separates arrival rate, service time, utilisation, and waiting time. However, it also contains a strong assumption: the service-time distribution is independent of load. A request that requires $S$ milliseconds of service at low load is assumed to require the same service-time distribution at higher load. This thesis shows that this assumption is not always valid for real containerised servers.

The project began as a teaching and capacity-planning platform. The original goal was to build a lightweight Docker-based laboratory where students could run microservices, generate open-loop load, collect request-level traces, and compare measured response-time distributions against discrete-event simulation (DES). As the platform matured, the work developed into a broader empirical investigation of when standard M/G/1 and M/G/$c$ models are adequate and when they fail. The key discovery was that several CPU-bound servers exhibit load-dependent wall-clock service time: as nominal utilisation increases, measured service time itself grows.

This finding changes the modelling problem. If service time changes with load, then a DES calibrated at low load will underestimate response times at higher load even if the arrival process, worker count, and base service-time distribution are correct. To quantify the effect, this thesis fits a family of service-time models:

\[
M0: E[S] = S_0,
\qquad
M1: E[S] = S_0 + \alpha\rho_0,
\qquad
M2: E[S] = \frac{S_0}{1-\beta\rho_0},
\]

where $S_0$ is the base service time and $\rho_0=\lambda S_0/c$ is nominal utilisation. The hyperbolic M2 model is motivated by processor-sharing-style slowdown, but is interpreted carefully as a reduced-form correction for scheduler-induced wall-clock service inflation rather than as a literal nested queue. An additional exploratory M3 contention-penalty form is also tested to check whether the same effect can be expressed as base service work plus a contention penalty.

The empirical platform covers multiple server implementations: Apache/PHP DSP, Apache message handling, Node.js DSP, Python/Gunicorn DSP, Java DSP, Go lognormal, and Go SQLite, with both single-core and multi-core variants where appropriate. These servers were deliberately chosen to expose different concurrency structures: OS process-per-request execution, event loops, Python GIL behaviour, Java JIT warmup, Go runtime scheduling, SQLite/WAL locking, and Node.js cluster splitting. This server catalogue allows the thesis to distinguish between three kinds of model failure: load-dependent service time, wrong queueing topology, and runtime measurement artefacts.

The thesis validates the models using response-time CDFs and Kolmogorov--Smirnov (KS) distance, not only mean response time. A service-time model is therefore considered useful only if it improves the simulated response-time distribution, not merely if it fits the measured mean service time. Supervisor follow-up work led to explicit CDF\_REF vs.\ CDF\_M0/M1/M2 comparisons, service-time histograms by arrival rate, and ANOVA tests on $\log(\texttt{service\_ms})$. These provide experimental, simulation, and statistical evidence for the same conclusion: the constant-service-time assumption fails for specific CPU-bound containerised workloads.

The work also studies boundaries. Node.js cluster is not well represented as a shared M/G/3 queue; it behaves closer to independent M/G/1 queues fed by connection splitting. Go SQLite points toward a tandem queue, where an application-level queue feeds an internal database/WAL lock queue. Java DSP is affected by JIT warmup, producing a runtime phase change rather than a queueing effect. Finally, exploratory priority/preemptive-resume and feedback-priority DES variants were implemented and validated. They improve over plain FCFS M0 in some selected traces, but do not outperform the simpler M2 load-dependent correction overall. This makes them promising future directions only if new two-class traces are collected.

The thesis therefore has two connected contributions. First, it delivers a reproducible Docker capacity-modelling laboratory suitable for teaching and experimentation. Second, it contributes a mechanistic, low-calibration-cost modelling workflow for identifying when service time is load-dependent and when a standard M/G/$c$ model is insufficient. The broader industrial relevance is proactive capacity planning: a future system can combine arrival-rate forecasting with the fitted service response model to predict SLO risk before queues build, rather than reacting after latency or CPU thresholds have already been breached.

\section{Thesis Structure}

Chapter~\ref{ch:background} introduces the queueing, DES, utilisation, and operational-analysis concepts needed to interpret the experiments. Chapter~\ref{ch:aims} states the research aims, objectives, questions, scope, and contributions. Chapter~\ref{ch:litreview} reviews related work in container performance, capacity planning, hybrid ML--DES modelling, MLASP, and educational virtualisation. Chapter~\ref{ch:methodology_b} describes the experimental platform, server catalogue, trace format, load generation, and DES validation method. Chapter~\ref{ch:journey} documents the research process and the debugging discoveries that shaped the final model. Chapter~\ref{ch:loaddep} presents the load-dependent service-time model, M0/M1/M2 validation, M3 stress test, histograms, ANOVA, and exploratory priority/feedback modelling. Chapter~\ref{ch:conclusion_b} summarises the contributions, limitations, and future directions.
